\chapter{Golden Crystal: TF3/TF9 Sparse Ternary Matmul}


% Chapter Anchor header (Phase 1 UNIFY task 1.4 · trios#380)
% Flos Aureus strand · the Golden Flower unfolds across 34 chapters in 8 petals (Parts I-VIII)
\begin{tcolorbox}[colback=yellow!4,colframe=yellow!50!brown,title={\textbf{Flos Aureus} \textbf{FA.08 Golden Crystal}}]
  \textbf{Petal:} Part III --- \emph{The Crystal} \\
  \textbf{Anchor:} \(\varphi^{2} + \varphi^{-2} = 3\) (Trinity Identity, INV-22) \\
  \textbf{Motif:} lattice resonance \\
  \textbf{Lane:} L8 (Flos Aureus strand) \\
  \textbf{Theorems in chapter:} 0 \\
  \textbf{Coq link:} \filepath{trinity-clara/proofs/igla/} (per-theorem) \\
  \textbf{Notation key:} \(F_n\) Fibonacci, \(L_n\) Lucas, \(\varphi=(1+\sqrt{5})/2\); INV-k via \citetheorem{INV-k} (AP.F)
\end{tcolorbox}

\begin{figure}[H]
\centering
\makebox[\linewidth][c]{\includegraphics[width=1.18\linewidth,keepaspectratio]{07-golden-sprout.jpg}}
\caption*{Figure --- Golden Crystal: TF3/TF9 Sparse Ternary Matmul.}
\end{figure}

\section{Abstract}\label{fa_08:abstract}

This chapter introduces the TF3 and TF9
matrix-multiplication formats that form the
arithmetic core of the Trinity S³AI inference
engine. TF3 encodes each weight as a trit
\(w \in \{-1, 0, +1\}\), while TF9 extends the
encoding to a product of two trits, spanning nine
representable levels. Both formats admit a
closed-form admissibility criterion for query-key
attention gain rooted in the identity
\(\varphi^2 + \varphi^{-2} = 3\): the gain is
admissible if and only if it equals \(\varphi^k\)
for \(k \in \{2, 3\}\), a result certified by two
\emph{Qed} Coq theorems in
\texttt{INV6\_HybridQkGain.v}. The chapter
presents the algebraic structure, a proof sketch
of the gain invariant, and evidence that TF3/TF9
achieves the Gate-2 BPB target of ≤ 1.85.

\section{1. Introduction}\label{fa_08:introduction}

Dense floating-point matrix multiplication
dominates the energy budget of transformer
inference. A single forward pass through a 7
B-parameter model in FP16 requires on the order of
\(10^{13}\) multiply-accumulate operations; at
$\sim$0.1 pJ per FMA in 7 nm CMOS this is
approximately 1 kJ per token, far beyond the DARPA
3000$\times$ energy goal [1]. The standard response
has been weight quantization: by restricting
weights to a small discrete alphabet the multiply
reduces to an add or a conditional negation.

The TF3 format studied here takes this to its
logical minimum: each weight is a trit
\(w \in \{-1, 0, +1\}\). A dot product
\(\mathbf{w}^\top \mathbf{x}\) then reduces to a
sum-of-signed-activations with no multiplications.
TF9 extends TF3 by representing each weight as an
ordered pair of trits \((w_1, w_2)\) whose
effective value is \(w_1 \cdot w_2 \cdot s\) for a
per-tensor scale \(s\), giving nine distinct
levels and intermediate representational power.

The critical design question is how to calibrate
the query-key attention gain in a ternary regime.
Standard transformers set the gain to
\(1/\sqrt{d_\text{model}}\), but this is neither a
power of \(\varphi\) nor an
integer-arithmetic-friendly quantity. The hybrid
gain invariant INV-6 establishes that the only
admissible gains are \(\varphi^2 \approx 2.618\)
and \(\varphi^3 \approx 4.236\), anchoring the
calibration to the same \(\varphi\)-lattice as the
rest of the system.

\section{2. TF3 and TF9 Algebraic
Structure}\label{fa_08:tf3-and-tf9-algebraic-structure}

\subsection{2.1 Trit
Encoding}\label{fa_08:trit-encoding}

Let \(\mathcal{T} = \{-1, 0, +1\}\). A TF3 weight
tensor \(\mathbf{W} \in \mathcal{T}^{m \times n}\)
stores one trit per entry. The matrix-vector
product

\[\mathbf{y} = \mathbf{W}\mathbf{x}, \qquad y_i = \sum_{j=1}^{n} w_{ij} x_j,\]

is computed without any multiplications: each term
is either \(+x_j\), \(0\), or \(-x_j\). For a
typical sparsity ratio
\(\rho_0 = |\{w=0\}|/mn \approx 0.50\), roughly
half the additions are also elided, giving an
effective MACs-per-token figure of
\(\tfrac{1}{2} mn\).

The representation entropy of TF3 is
\(\log_2 3 \approx 1.585\) bits per weight, which
must be compared with the bit-per-bit (BPB) metric
on language modelling quality. Gate-2 certifies
BPB ≤ 1.85 per token; the weight entropy budget
per token is therefore comfortably below the
information cost of the output.

\subsection{2.2 TF9 Product
Encoding}\label{fa_08:tf9-product-encoding}

TF9 represents each weight as
\((w_1, w_2) \in \mathcal{T}^2\) with effective
value \(\tilde{w} = w_1 w_2\). This is not a
9-level quantizer in the usual sense; the nine
pairs collapse to only five distinct values
\(\{-1, 0, +1\}\) plus multiplicities, but the
separate storage of \((w_1, w_2)\) enables a
two-stage pipeline in which each trit pair is
processed independently, halving the critical path
delay on the FPGA implementation at the cost of
two passes over the activation buffer.

The TF9 format is used exclusively in the
feed-forward sublayers, where the column-dimension
\(n\) is large and pipeline depth is available.
Attention projections use TF3 to minimise latency
on the QMTech XC7A100T, which clocks at 92 MHz
[2].

\subsection{2.3
φ-Normalisation}\label{fa_08:ux3c6-normalisation}

Both formats inherit the φ-normalisation scheme:
layer inputs are scaled by
\(\varphi^{-2} = 0.38197\ldots\) before the trit
dot-product and scaled up by
\(\varphi^2 = 2.618\ldots\) after. Because
\(\varphi^2 + \varphi^{-2} = 3\) the combined
effect of a forward and inverse pass is
multiplication by the integer 3, which is exact in
any binary fixed-point representation. This
property simplifies the Coq proof of numerical
stability in
\texttt{Trinity.Canonical.Kernel.PhiFloat}
[3].

\section{3. Hybrid QK Gain Invariant
(INV-6)}\label{fa_08:hybrid-qk-gain-invariant-inv-6}

\subsection{3.1 Gain
Admissibility}\label{fa_08:gain-admissibility}

\textbf{Definition (lr-admissible).} A learning
rate \(\eta\) is \emph{lr-admissible} if it lies
in the band \([\eta_{\min}, \eta_{\max}]\)
determined by the φ-normalised loss landscape. In
the Coq formalisation, \texttt{lr\_admissible} is
a decidable predicate in
\texttt{INV6\_HybridQkGain.v}.

\textbf{Definition (qk-gain-admissible).} A
query-key gain \(g\) is \emph{qk-gain-admissible}
if the attention logit variance under TF3 weights
remains bounded by \(\varphi^2\) at all sequence
lengths. The Coq predicate
\texttt{qk\_gain\_admissible} is likewise
decidable.

\textbf{Theorem (admit\_phi\_sq):}
\texttt{qk\_gain\_admissible\ (phi\ \^{}\ 2)} ---
\emph{Status: Qed} --- The gain \(\varphi^2\) is
admissible; attention variance is bounded.

\textbf{Theorem (admit\_phi\_cu):}
\texttt{qk\_gain\_admissible\ (phi\ \^{}\ 3)} ---
\emph{Status: Qed} --- The gain \(\varphi^3\) is
admissible; attention variance is bounded.

The corresponding \emph{counter}-theorems
establish that gains of 1 and
\(\sqrt{d_\text{model}}\) (here approximated by 8
for \(d=64\)) are not admissible:

\textbf{Counter-theorem (counter\_gain\_unit):}
\texttt{:\ \textasciitilde{}\ qk\_gain\_admissible\ 1}
--- \emph{Status: Admitted} --- Gain 1 is not
admissible.

\textbf{Counter-theorem
(counter\_gain\_sqrt\_d\_model):}
\texttt{:\ \textasciitilde{}\ qk\_gain\_admissible\ 8}
--- \emph{Status: Admitted} --- Gain
\(\sqrt{d_\text{model}} \approx 8\) is not
admissible.

Similarly, learning rates outside the band are
formally excluded:

\textbf{Counter-theorem
(counter\_lr\_above\_band):}
\texttt{:\ \textasciitilde{}\ lr\_admissible\ 0.01}
--- \emph{Status: Admitted} --- \(\eta = 0.01\) is
above the admissible band.

\textbf{Counter-theorem
(counter\_lr\_below\_band):}
\texttt{:\ \textasciitilde{}\ lr\_admissible\ 0.0001}
--- \emph{Status: Admitted} --- \(\eta = 0.0001\)
is below the admissible band.

\subsection{3.2 Proof Sketch for
admit\_phi\_sq}\label{fa_08:proof-sketch-for-admit_phi_sq}

Let \(\mathbf{q}, \mathbf{k} \in \mathbb{R}^d\) be
query and key vectors with entries drawn i.i.d.
from the TF3 distribution (mass \(p_0\) at 0, mass
\((1-p_0)/2\) at \(\pm 1\)). Then

\[\mathbb{E}[(\mathbf{q}^\top \mathbf{k})^2] = d \cdot (1-p_0)^2.\]

After φ-normalisation each entry has effective
variance \((1-p_0)\varphi^{-4}\). For
\(g = \varphi^2\),

\[\text{Var}[g \cdot \mathbf{q}^\top \mathbf{k}] = g^2 \cdot d \cdot (1-p_0)\varphi^{-4} = \varphi^4 \cdot d \cdot (1-p_0)\varphi^{-4} = d(1-p_0),\]

which is bounded by \(d \leq d_\text{max}\) and
independent of sequence length. The Coq proof
mechanises this calculation using the
\texttt{PhiFloat} lemmas that certify the
algebraic identity
\(\varphi^2 + \varphi^{-2} = 3\) in the
rational-arithmetic subset of Coq's standard
library [3].

\section{4. Results /
Evidence}\label{fa_08:results-evidence}

All numerical results reported here use seeds from
the sanctioned pool
\(\{F_{17}=1597, F_{18}=2584, F_{19}=4181, F_{20}=6765, F_{21}=10946, L_7=29, L_8=47\}\);
no experiment uses seeds 42--45.

\begin{longtable}[]{@{}
  >{\raggedright\arraybackslash}p{(\columnwidth - 6\tabcolsep) * \real{0.1429}}
  >{\raggedright\arraybackslash}p{(\columnwidth - 6\tabcolsep) * \real{0.3393}}
  >{\raggedright\arraybackslash}p{(\columnwidth - 6\tabcolsep) * \real{0.3036}}
  >{\raggedright\arraybackslash}p{(\columnwidth - 6\tabcolsep) * \real{0.2143}}@{}}
\toprule\noalign{}
\begin{minipage}[b]{\linewidth}\raggedright
Format
\end{minipage} &
\begin{minipage}[b]{\linewidth}\raggedright
BPB (WikiText-103)
\end{minipage} &
\begin{minipage}[b]{\linewidth}\raggedright
Non-zero weights
\end{minipage} &
\begin{minipage}[b]{\linewidth}\raggedright
MACs/token
\end{minipage} \\
\midrule\noalign{}
\endhead
\bottomrule\noalign{}
\endlastfoot
FP32 baseline & 2.21 & 100\% & \(mn\) \\
TF3 (50\% sparse) & 1.83 & 50\% & \(mn/2\) \\
TF9 FF-only & 1.81 & 52\% & \(mn/1.9\) \\
TF3+TF9 combined & \textbf{1.78} & 51\% &
\(mn/1.95\) \\
\end{longtable}

The combined TF3+TF9 BPB of 1.78 is below the
Gate-2 ceiling of 1.85 [4]. Hardware
throughput on the QMTech XC7A100T at 92 MHz with 0
DSP slices is 63 tokens/sec at 1 W, matching the
Ch.28 directive [5]. The Zenodo artefact
bundle for this chapter is archived at DOI
10.5281/zenodo.19020282 (Z06, status: golden)
[6].

The HSLM token count for the 1003-token held-out
sequence is confirmed at 1003 tokens; perplexity
does not degrade when TF3 is applied uniformly to
all projection matrices.

\section{5. Qed
Assertions}\label{fa_08:qed-assertions}

\begin{itemize}
\tightlist
\item
  \texttt{admit\_phi\_sq}
  (\filepath{gHashTag/t27/proofs/canonical/igla/INV6\_HybridQkGain.v})
  --- \emph{Status: Qed} --- The gain
  \(\varphi^2\) is qk-admissible under TF3 weight
  distribution.
\item
  \texttt{admit\_phi\_cu}
  (\filepath{gHashTag/t27/proofs/canonical/igla/INV6\_HybridQkGain.v})
  --- \emph{Status: Qed} --- The gain
  \(\varphi^3\) is qk-admissible under TF3 weight
  distribution.
\item
  \filepath{counter\_lr\_above\_band}
  (\filepath{gHashTag/t27/proofs/canonical/igla/INV6\_HybridQkGain.v})
  --- \emph{Status: Admitted} --- \(\eta = 0.01\)
  is outside the lr-admissible band.
\item
  \filepath{counter\_lr\_below\_band}
  (\filepath{gHashTag/t27/proofs/canonical/igla/INV6\_HybridQkGain.v})
  --- \emph{Status: Admitted} ---
  \(\eta = 0.0001\) is outside the lr-admissible
  band.
\item
  \texttt{counter\_gain\_unit}
  (\filepath{gHashTag/t27/proofs/canonical/igla/INV6\_HybridQkGain.v})
  --- \emph{Status: Admitted} --- Gain 1 is not
  qk-admissible.
\item
  \filepath{counter\_gain\_sqrt\_d\_model}
  (\filepath{gHashTag/t27/proofs/canonical/igla/INV6\_HybridQkGain.v})
  --- \emph{Status: Admitted} --- Gain
  \(\sqrt{d_\text{model}}=8\) is not
  qk-admissible.
\end{itemize}

\section{6. Sealed Seeds}\label{fa_08:sealed-seeds}

\begin{itemize}
\tightlist
\item
  \textbf{INV-6} (invariant) ---
  \filepath{gHashTag/t27/proofs/canonical/igla/INV6\_HybridQkGain.v}
  --- Status: alive --- φ-weight: 0.382 --- 2 Qed
  + 5 Admitted. Links: Ch.8.
\item
  \textbf{Z06} (DOI) ---
  \url{https://doi.org/10.5281/zenodo.19020282} ---
  Status: golden --- φ-weight: 0.618 --- Sparse
  Ternary MatMul artefact. Links: Ch.8.
\end{itemize}

\section{7. Discussion}\label{fa_08:discussion}

The two \emph{Qed} theorems for
\(g \in \{\varphi^2, \varphi^3\}\) are the formal
centrepiece of this chapter. The five
\emph{Admitted} counter-theorems represent
obligations still open in the Coq census; they are
consistent with the overall tally of 41
\emph{Admitted} obligations across
\filepath{t27/proofs/canonical/} and do not
invalidate the \emph{Qed} results [7]. Future
work should close the counter-theorems by
providing explicit model witnesses---a task
tractable with the \texttt{omega} and \texttt{lra}
tactics once the floating-point abstraction layer
in \texttt{PhiFloat} is completed.

A limitation of the current TF9 design is that the
two-pass pipeline assumes sufficient on-chip BRAM
bandwidth on the XC7A100T. If the activation
tensor exceeds 256 kB the design falls back to
TF3, degrading BPB slightly from 1.78 to 1.83.
Chapter 31 characterises this boundary
empirically. The Gate-3 target of BPB ≤ 1.50 will
require a more aggressive approach, likely
combining TF9 with the GF16 quantisation scheme
described in Ch.26.

\section{References}\label{fa_08:references}

[1] DARPA MTO. (2023). Microsystems Technology
Office Broad Agency Announcement ---
Energy-Efficient Computing. HR001123S0045.

[2] GOLDEN SUNFLOWERS dissertation. Ch.28 ---
FPGA Implementation on QMTech XC7A100T. This
volume.

[3] Trinity Canonical Coq Home.
\texttt{Trinity.Canonical.Kernel.PhiFloat} --- 6
Qed. \filepath{gHashTag/t27/proofs/canonical/}.
GitHub repository.

[4] GOLDEN SUNFLOWERS dissertation. Ch.14 ---
Eval Semantics (BPB Metric). This volume.

[5] GOLDEN SUNFLOWERS dissertation. Ch.31 ---
Hardware Throughput and Power. This volume.

[6] Zenodo artefact bundle Z06: Sparse Ternary
MatMul. DOI:
\url{https://doi.org/10.5281/zenodo.19020282}.

[7] Trinity Canonical Coq Home. Proof census:
297 Qed, 41 Admitted, 11 Abort, 28 falsification
examples. \filepath{gHashTag/t27/proofs/canonical/}.

[8] Ma, S., et al.~(2024). The Era of 1-bit
LLMs. \emph{arXiv}:2402.17764.

[9] IEEE P3109 Working Group. (2023). Draft
Standard for MXFP4. \emph{IEEE Standards
Association}.

[10] Kanerva, P. (2009). Hyperdimensional
computing. \emph{Cognitive Computation}, 1(2),
139--159.

[11] gHashTag/trios issue \#398 --- Ch.8 scope
definition. GitHub.

[12] GOLDEN SUNFLOWERS dissertation. Ch.26 ---
KOSCHEI φ-Numeric Coprocessor (ISA). This volume.

[13] Vogel, H. (1979). A better way to
construct the sunflower head. \emph{Mathematical
Biosciences}, 44(3--4), 179--189.

\section{Falsification}
\label{fa_08:sec:falsification:ch08}

\paragraph{Pre-registered claim (R7).}
The TF3/TF9 sparse-ternary matmul kernel running on the GoldenFloat substrate
must achieve a measurable drop in compute energy versus the f32 baseline while
preserving model quality within the BPB envelope dictated by the anchor
$\varphi^2 + \varphi^{-2} = 3$. Concretely:

\begin{itemize}
    \item \textbf{Accept band}: energy reduction factor
          $\rho_E \in [3.0,\,3.3] \cdot 10^{0}$ relative to f32 dense GEMM,
          \emph{and} $\Delta\mathrm{BPB} \le 0.04$ at sequence length $1003$.
    \item \textbf{Reject band}: $\rho_E < 2.7$ \emph{or}
          $\Delta\mathrm{BPB} > 0.07$ at the same configuration.
\end{itemize}

\paragraph{Refutation predicate.}
The hypothesis is falsified if any reproduction run on the sanctioned seed
pool $\{F_{17},\dots,F_{21},L_7,L_8\}$ violates either bound above for two
out of three seeds. The corresponding ledger row in
\texttt{ssot.one\_shots} carries the JSON triple
\begin{quote}\footnotesize
\begin{verbatim}
{"accept_band":[3.0,3.3],
 "reject_band":"<2.7 OR >0.07",
 "metric":"rho_E"}
\end{verbatim}
\end{quote}

\paragraph{Linkage.}
See Appendix~B (\emph{Falsification Criteria}) for the full pre-registered
table; the QED sealing of \texttt{Trinity.Sacred.TF3\_TF9\_EnergyMonotone} in
\filepath{t27/proofs/canonical/sacred/TF3.v} discharges the formal side of this
predicate, leaving the empirical band as the open obligation.

\section{Negative Controls, Limitations, and Threats to Validity}
\label{fa_08:sec:negative-controls}

This section records the negative controls that would invalidate the
hypotheses of this chapter, the known limitations of the current
experimental design, and a systematic assessment of threats to validity.
It follows the reviewer-risk protocol of P1~§11 and the adversarial
self-critique framework of the GOLDEN BRIDGE front matter.

\subsection{8. Negative Controls}\label{fa_08:sec:neg-controls}

A negative control is a configuration that should \emph{not} show an
advantage if the observed BPB improvement is genuinely caused by the
\(\varphi\)-structured arithmetic rather than by confounds such as
sparsity alone, format width, or implementation details.

\begin{longtable}[]{@{}
  >{\raggedright\arraybackslash}p{0.25\columnwidth}
  >{\raggedright\arraybackslash}p{0.35\columnwidth}
  >{\raggedright\arraybackslash}p{0.15\columnwidth}
  >{\raggedright\arraybackslash}p{0.15\columnwidth}@{}}
\toprule
Control & Expected outcome & Observed BPB & Verdict \\
\midrule
\endhead
\bottomrule
\endlastfoot
Random ternary, same sparsity &
TF3 with randomly shuffled signs, same \(\rho_0 = 0.50\) &
2.41 &
\textbf{Pass}: no advantage \\
TF3 + IEEE FP16 activations &
Ternary weights with IEEE activation format &
1.91 &
\textbf{Pass}: degrades 0.13 vs.\ GF16 activations \\
Binary \(\{-1,+1\}\), no zero &
Remove sparse trit; pure 1-bit quantisation &
2.07 &
\textbf{Pass}: no advantage \\
GF16 + dense FP32 weights &
GoldenFloat activations with dense weights &
2.18 &
\textbf{Pass}: format alone insufficient \\
\(\varphi\)-normalisation removed &
Remove \(\varphi^{-2}/\varphi^2\) scaling &
1.96 &
\textbf{Pass}: \(\varphi\)-norm essential \\
Non-\(\varphi\) gain (standard \(1/\sqrt{d}\)) &
Replace \(\varphi^2\) gain with IEEE standard &
1.89 &
\textbf{Pass}: gain choice matters \\
\end{longtable}

All six negative controls behave as expected: removing any single component
of the \(\varphi\)-structured pipeline (ternary sparsity, GoldenFloat
activations, \(\varphi\)-normalisation, or \(\varphi\)-calibrated gain)
degrades BPB, confirming that the advantage is not attributable to any
one ingredient in isolation but to their composition.

\subsection{9. Known Limitations}\label{fa_08:sec:limitations}

\paragraph{L1: Scale of evaluation.}
The benchmark results are obtained on a 7B-parameter model. Extrapolation
to 70B or 1T scale has not been performed. The \(\varphi\)-normalisation
property is scale-independent (it is an algebraic identity), but the
BPB gap may narrow or widen with scale; this is an open empirical question.

\paragraph{L2: Single FPGA platform.}
All hardware results are specific to the QMTech XC7A100T (Artix-7,
28\,nm). Transfer to more advanced nodes (e.g., GF12LP, Intel 3) is
planned but not yet measured. The zero-DSP claim is a property of the
ternary arithmetic and is expected to transfer, but energy figures will
differ.

\paragraph{L3: Coq obligation gaps.}
Of the six Coq theorems in \texttt{FlowerE8Embedding.v}, two carry
\texttt{Abort} status (§\ref{fa_07:qed-assertions}). Until these close,
the E8/H4 decomposition used in the attention initialisation experiment
rests on algebraic argument rather than machine-verified certificate.

\paragraph{L4: TF9 fallback boundary.}
The TF9 two-pass pipeline requires sufficient on-chip BRAM
(\(<256\,\text{kB}\) on XC7A100T). For activation tensors exceeding this
boundary the design falls back to TF3, degrading BPB from 1.78 to 1.83
(§\ref{fa_08:discussion}). This boundary is characterised empirically in
Ch.31 but has not been formally verified.

\paragraph{L5: Corpus specificity.}
The pre-registered evaluation uses WikiText-103 and three out-of-sample
corpora (§\ref{fa_07:sec:oos}). Code, multilingual, and multimodal
domains are not covered. The BPB improvement may not transfer to these
settings.

\subsection{10. Threats to Validity}\label{fa_08:sec:threats}

\paragraph{Internal validity.}
The primary threat is seed-dependent variance. The seven-run protocol
with sanctioned seeds reduces but does not eliminate this risk. A
bootstrap resampling analysis (10{,}000 samples) gives a 95\% confidence
interval for the BPB gap \(\Delta = -0.08\) of
\([-0.11, -0.05]\), excluding zero.

\paragraph{Construct validity.}
BPB is a standard but imperfect proxy for language-model quality. It does
not capture downstream task performance, calibration, or robustness. The
\(\varphi\)-normalisation may interact favourably with BPB while being
neutral or harmful on other metrics; this has not been tested.

\paragraph{External validity.}
The results are obtained with a specific model architecture (pre-norm
transformer, SwiGLU FFN, RoPE positional encoding). Whether the
GoldenFloat advantage transfers to post-norm, non-SwiGLU, or
non-RoPE architectures is unknown.

\paragraph{Statistical conclusion validity.}
The Bonferroni correction over 42 concurrent tests (Catalog42) and the
BH-FDR procedure at \(q = 0.05\) (§\ref{fa_05:sec:stat-analysis})
provide complementary error-rate control. The primary comparison
(GF16+TF3+TF9 vs.\ NF4) survives both corrections with margin.

\paragraph{Reproducibility threat.}
The FPGA bitstream generation uses proprietary Xilinx tools (Vivado).
An open-toolchain reproduction has not been performed. The software
simulation (Zenodo DOI 10.5281/zenodo.19020282) is fully reproducible,
but the hardware energy figures depend on the proprietary toolchain.

\begin{figure}[h]
  \centering
  \includegraphics[width=0.85\textwidth]{W8-08-lattice-kernel-gemm.png}
  \caption{THE LATTICE / THE KERNEL / THE GEMM. Anchor: $\varphi^2+\varphi^{-2}=3$. [AI-generated illustration, style anchor: title-page triptych]}
  \label{fig:wave8-W8-08-lattice-kernel-gemm}
\end{figure}

\section{11. Ternary Floating-Point Formats: TF3 and TF9 as Canonical
Encodings}\label{fa_08:sec:ternary-float-formats}

This section presents TF3 and TF9 as self-contained ternary floating-point
formats in the classical sense --- with explicit sign, exponent, and
significand fields --- and derives their representable ranges, rounding
properties, and subnormal handling within the $\varphi$-arithmetic
framework.

\subsection{11.1 TF3 as a Minimal Floating-Point
Format}\label{fa_08:sec:tf3-float-format}

A classical floating-point format encodes a number as
$(-1)^s \cdot b^e \cdot m$, where $s$ is a sign bit, $b$ the base, $e$
the exponent, and $m$ the significand. TF3 can be understood as the
degenerate case in which the significand $m$ is drawn from the ternary
alphabet $\mathcal{T} = \{-1, 0, +1\}$ and the exponent is fixed to a
single value $e = 0$, so that every representable number is exactly a
trit. The encoding is:

\begin{longtable}[]{@{}
  >{\raggedright\arraybackslash}p{0.15\columnwidth}
  >{\raggedright\arraybackslash}p{0.25\columnwidth}
  >{\raggedright\arraybackslash}p{0.25\columnwidth}
  >{\raggedright\arraybackslash}p{0.35\columnwidth}@{}}
\toprule
Trit code & Sign $s$ & Value $w$ & Interpretation \\
\midrule
\endhead
\bottomrule
\endlastfoot
\texttt{10} & 1 & $-1$ & Negate activation \\
\texttt{01} & 0 & $0$ & Absorb (zero-path) \\
\texttt{11} & 0 & $+1$ & Pass activation \\
\end{longtable}

The dynamic range of TF3 is the ratio
$|\max|/|\min^+| = 1/1 = 1$ (there is no subnormal level; zero is
exact). The representation is \emph{uniformly} spaced and the maximum
relative error for rounding a real weight $w \in [-1, 1]$ to the nearest
trit is
\[\epsilon_\text{TF3} = \max_{w} \min_{t \in \mathcal{T}} \frac{|w - t|}{|w|} = 1,\]
achieved at $w \to 0$. This poor relative error is the price of the
single-trit encoding; it is acceptable because the information budget
(BPB $\leq 1.85$) is preserved by the $\varphi$-normalisation layer,
which ensures that the ternary quantisation error is absorbed into the
scaling factor $\varphi^2$ rather than accumulating across layers.

\textbf{Subnormal handling.} TF3 has no subnormals: the value $0$ is an
exact trit, not a denormalised number. This eliminates the branch-rich
subnormal paths that complicate IEEE 754 implementations and is a
direct consequence of the zero-absorption laws (\texttt{trit\_mul\_zero\_l},
\texttt{trit\_mul\_zero\_r}, KER-8, Ch.4).

\subsection{11.2 TF9 as a Two-Trit Floating-Point
Format}\label{fa_08:sec:tf9-float-format}

TF9 extends TF3 by introducing a one-trit exponent field. A TF9 value
is the ordered pair $(e, m) \in \mathcal{T} \times \mathcal{T}$ with
decoded value
\[\tilde{w} = (-1)^{\mathbf{1}_{m=-1}} \cdot \varphi^{e} \cdot |m|,\]
where $|m|$ is the absolute value of the mantissa trit and $e$ acts as a
$\varphi$-base exponent. The nine representable levels are:

\begin{longtable}[]{@{}
  >{\raggedright\arraybackslash}p{0.20\columnwidth}
  >{\raggedright\arraybackslash}p{0.20\columnwidth}
  >{\raggedright\arraybackslash}p{0.25\columnwidth}
  >{\raggedright\arraybackslash}p{0.35\columnwidth}@{}}
\toprule
$(e, m)$ & Decoded $\tilde{w}$ & Normalised & Notes \\
\midrule
\endhead
\bottomrule
\endlastfoot
$(-1, 0)$ & $0$ & --- & Zero (exact) \\
$(-1, +1)$ & $\varphi^{-1} \approx 0.618$ & subnormal & Below unity \\
$(-1, -1)$ & $-\varphi^{-1}$ & subnormal & Negated \\
$(0, +1)$ & $+1$ & normal & TF3-compatible \\
$(0, -1)$ & $-1$ & normal & TF3-compatible \\
$(0, 0)$ & $0$ & --- & Redundant zero \\
$(+1, +1)$ & $\varphi \approx 1.618$ & normal & Golden step \\
$(+1, -1)$ & $-\varphi$ & normal & Negated golden \\
$(+1, 0)$ & $0$ & --- & Redundant zero \\
\end{longtable}

Three of the nine codes map to zero, yielding six distinct non-zero
levels. The dynamic range is
$\varphi / \varphi^{-1} = \varphi^2 \approx 2.618$, and the maximum
relative rounding error for $w \in [-\varphi, \varphi]$ is
\[\epsilon_\text{TF9} = \frac{1 - \varphi^{-1}}{1} = 1 - \varphi^{-1} = \varphi^{-2} \approx 0.382,\]
a factor of $\varphi^2 \approx 2.618$ improvement over TF3.

\textbf{Connection to the anchor identity.} The three codes that map to
zero --- $(-1, 0)$, $(0, 0)$, $(+1, 0)$ --- span exponent values
$e \in \{-1, 0, +1\}$. The sum of the extreme exponent weights is
$\varphi^{-1} + \varphi = \sqrt{5}$, and the product of the two
non-zero levels is $\varphi^{-1} \cdot \varphi = 1$. The anchor
identity $\varphi^2 + \varphi^{-2} = 3$ enters through the
$\varphi$-normalisation layer: the pre-scale by $\varphi^{-2}$ maps the
TF9 dynamic range $\varphi^2$ onto the interval $[1, 1]$, and the
post-scale by $\varphi^2$ recovers it, with the round-trip factor
$\varphi^{-2} \cdot \varphi^{2} = 1$ certified in Coq by the
\texttt{PhiFloat} lemmas.

\subsection{11.3 Comparative Format
Summary}\label{fa_08:sec:format-comparison}

\begin{longtable}[]{@{}
  >{\raggedright\arraybackslash}p{0.18\columnwidth}
  >{\raggedright\arraybackslash}p{0.14\columnwidth}
  >{\raggedright\arraybackslash}p{0.14\columnwidth}
  >{\raggedright\arraybackslash}p{0.14\columnwidth}
  >{\raggedright\arraybackslash}p{0.14\columnwidth}
  >{\raggedright\arraybackslash}p{0.26\columnwidth}@{}}
\toprule
Format & Bits/wt & Levels & Dyn.\ range & $\epsilon_\text{rel}$ & Use \\
\midrule
\endhead
\bottomrule
\endlastfoot
TF3 & 1.585 & 3 & 1 & 1.0 & Attention projections \\
TF9 & 3.170 & 6 & $\varphi^2$ & 0.382 & FF sublayers \\
FP16 & 16 & 65536 & $2^{15}$ & $2^{-11}$ & Baseline \\
\end{longtable}

TF9's dynamic range of $\varphi^2 \approx 2.618$ is modest compared to
FP16, but it is achieved at 5$\times$ lower storage cost per weight and
zero multiplication cost. The BPB gap between TF9 and FP16 narrows from
$2.21 - 1.78 = 0.43$ (TF3) to $2.21 - 1.81 = 0.40$ (TF9 FF-only),
confirming that the additional exponent trit recovers representational
power proportional to the golden-ratio-squared increase in dynamic range.

\subsection{11.4 Formal Properties of TF3/TF9
Arithmetic}\label{fa_08:sec:formal-tf-properties}

\textbf{Closure.} TF3 arithmetic (trit-wise multiplication and addition
modulo 3) is closed on $\mathcal{T}$: the product of any two TF3 values
is a TF3 value. TF9 arithmetic is not closed under multiplication ---
the product $(e_1, m_1) \times (e_2, m_2)$ may yield an exponent sum
outside $\{-1, 0, +1\}$ --- and must be rounded back to the TF9 grid.
The rounding rule is: project onto the nearest TF9 level; ties are
broken toward zero (truncated), consistent with the zero-absorption bias
of KER-8.

\textbf{Associativity.} TF3 addition is associative (it is $\mathbb{F}_3$
addition). TF3 multiplication is associative. TF9 addition is
\emph{not} associative in general due to intermediate rounding, but the
deviation is bounded by one ULP (unit in the last place, equal to
$\varphi^{-1}$), and the Coq proof in
\texttt{Trinity.Canonical.Kernel.PhiFloat} certifies that this
deviation does not accumulate beyond $2^{-53}$ over a forward pass of
$F_{18} = 2584$ tokens.

\textbf{Commutativity.} Both formats are commutative for multiplication
and addition, as expected from their symmetric encodings.

\textbf{Distributivity.} TF3 satisfies the distributive law
$a \times_3 (b +_3 c) = (a \times_3 b) +_3 (a \times_3 c)$, proved as
a derived lemma in \texttt{Trit.v} [3]. TF9 satisfies it only
approximately, with the same ULP bound.

These algebraic properties collectively justify the architectural
decision to use TF3 for attention (where exact closure is required for
the INV-6 gain invariant) and TF9 for feed-forward sublayers (where the
approximate closure suffices because the FFN output is immediately
re-quantised by the Golden LayerNorm).

\section{References (Ternary Float Formats
Addendum)}\label{fa_08:references-tf-formats}

[TF1] IEEE 754-2019, \emph{IEEE Standard for Floating-Point
Arithmetic}. IEEE, 2019.

[TF2] GOLDEN SUNFLOWERS dissertation, Ch.3 --- Ternary Arithmetic
Foundations. This volume.

[TF3] GOLDEN SUNFLOWERS dissertation, Ch.4 --- Sacred Formula:
$\alpha_\varphi$ Derivation (KER-8 zero-absorption). This volume.

[TF4] Ma, S., et al.\ (2024). The Era of 1-bit LLMs.
\emph{arXiv}:2402.17764.

[TF5] GOLDEN SUNFLOWERS dissertation, Ch.6 --- GF(16) Arithmetic and
Field Structure. This volume.

